\section{Symmetry, and why the objection to it is fair \hfill \open{}}

\textbf{The numbers first, because they are large.} The benchmark models are $N$ copies of one
template. Permuting the copies maps the system onto itself, so only one representative per orbit
needs exploring. For \texttt{fischer} with the query fixed, the number of orbits is exactly
\[
  12N^2 - 13N - 4 ,
\]
against a reachable state space that grows exponentially. That is $660$ orbits at $N=8$
($38\times$ fewer states), $1\,066$ at $N=10$ ($245\times$), and about $1\,568$ at $N=12$
($\approx 1750\times$). Exponential space, quadratic quotient.

\paragraph{The premise I used to dismiss it was wrong.} I rejected symmetry on the grounds that
``\uppaal{} already has it''. \uppaal{} has it \emph{for models written with \texttt{scalar} types}.
Our generators emit flattened templates with no scalar declarations, so \verifyta{} applies exactly
zero symmetry reduction to these inputs. Comparing our auto-detected symmetry against their
un-annotated model is therefore not a like-for-like comparison in our favour --- it is a comparison
of two different modelling styles.

\paragraph{And the objection to it is still fair.} The concern raised was that these benchmarks are
\emph{highly symmetric by construction}. Every one of them is ``$N$ identical replicas plus a shared
resource''. A technique whose payoff is proportional to how replicated the model is will look
spectacular here and do nothing on an industrial model where the components differ. That is a
benchmark-shape argument, not a soundness argument, and it is correct.

\paragraph{So what would make it honest.} Two things, and they are cheap:
\begin{enumerate}
\item Report symmetry results \emph{only} against a \uppaal{} model that also declares
      \texttt{scalar} types, so both tools reduce. Anything else is comparing annotations.
\item Include at least one asymmetric model, so the reader can see what the technique does when
      the replication is not there.
\end{enumerate}

\paragraph{Verdict.} Open, but the burden is on presentation rather than on the technique. The
$245\times$ is real; it is just measured on the shape most flattering to it. My earlier
justification for dismissing it was factually wrong, and the honest reason to be cautious is the
one that was raised: the benchmarks are too symmetric to prove anything general.
